\chapter{Messprotokolle und Rohdaten der Störungsexperimente}
\appendixentry{Messprotokolle und Rohdaten der Störungsexperimente}
\label{app:messprotokolle}

In diesem Anhang sind die maschinenlesbaren Konfigurationen, Ausführungsbefehle und Rohdatenprotokolle der drei durchgeführten Störungsexperimente dokumentiert.

\section{Experiment 1: Prozessabsturz und Redundanzvergleich}
\label{app:exp1-daten}
Zur Durchführung des Last- und Störungstests wurde das nachfolgende Python-Kommando ausgeführt. Es sendet kontinuierlich 15 Anfragen pro Sekunde und induziert bei Sekunde 12 über den Labor-Endpunkt einen Prozessabsturz:
\begin{lstlisting}[language=bash, caption={Ausführungsbefehl zu Experiment 1}]
python experiments/exp1_crash_test.py \
  --url http://63.176.27.134/api/info \
  --crash-url http://63.176.27.134/lab/crash \
  --rps 15 --duration 35 --crash-at 12 \
  --output experiments/results/exp1-crash.json
\end{lstlisting}

\begin{lstlisting}[language=json, caption={Rohdatenprotokoll exp1-crash.json}, basicstyle=\fontsize{8}{10}\ttfamily]
{
  "experiment": "Experiment 1 - Prozessabsturz (1 Pod vs 2 Pods)",
  "timestamp_iso": "2026-10-01T08:03:02.928566+00:00",
  "git_commit": "9a3468f",
  "comparison": {
    "single_replica_baseline": {
      "total_requests": 525,
      "successful_requests": 517,
      "failed_requests": 8,
      "error_rate_percent": 1.52,
      "latency_p50_ms": 50.75,
      "latency_p95_ms": 162.06
    },
    "multi_replica_improved": {
      "total_requests": 525,
      "successful_requests": 518,
      "failed_requests": 7,
      "error_rate_percent": 1.33,
      "latency_p50_ms": 55.2,
      "latency_p95_ms": 159.81
    }
  }
}
\end{lstlisting}

\section{Experiment 2: Fehlerhaftes Release und Git-Revert}
\label{app:exp2-daten}
Das nachfolgende Testskript simulierte ein fehlerhaftes Release mit fehlschlagender Readiness-Sonde und maß die Latenzen sowie die Wiederherstellungsdauer nach einem automatisierten \texttt{git revert}:
\begin{lstlisting}[language=bash, caption={Ausführungsbefehl zu Experiment 2}]
python experiments/exp2_bad_deploy.py \
  --repo-dir . \
  --service-url http://63.176.27.134/api/info \
  --output experiments/results/exp2-bad-deploy.json
\end{lstlisting}

\begin{lstlisting}[language=json, caption={Rohdatenprotokoll exp2-bad-deploy.json}, basicstyle=\fontsize{8}{10}\ttfamily]
{
  "experiment": "Experiment 2 - Fehlerhaftes Deployment und Git Revert",
  "timestamp_iso": "2026-10-01T08:04:58.168213+00:00",
  "commits": {
    "initial": "9a3468f11285c955a17e6980c57aeeef63274a73",
    "revert": "2d34aca6a4759d14acaf2d2af6c43ff1145ed189"
  },
  "traffic_during_bad_rollout": {
    "total_requests": 200,
    "successful_requests": 200,
    "failed_requests": 0,
    "error_rate_percent": 0.0
  },
  "recovery_time_seconds": 37.19
}
\end{lstlisting}

\section{Experiment 3: Konfigurationsdrift und Argo-CD-Self-Heal}
\label{app:exp3-daten}
Das Überwachungsskript skalierte das Deployment imperativ auf 0 Replikate und stoppte die Zeitspanne bis zur automatischen Selbstreparatur durch Argo~CD:
\begin{lstlisting}[language=bash, caption={Ausführungsbefehl zu Experiment 3}]
python experiments/exp3_config_drift.py \
  --deployment resilience-lab \
  --namespace default \
  --app-name resilience-lab \
  --output experiments/results/exp3-config-drift.json
\end{lstlisting}

\begin{lstlisting}[language=json, caption={Rohdatenprotokoll exp3-config-drift.json}, basicstyle=\fontsize{8}{10}\ttfamily]
{
  "experiment": "Experiment 3 - Konfigurationsabweichung und Self-Heal",
  "git_commit": "07a858dac10fe5e76bec13440a743e7a9331bd21",
  "metrics": {
    "detection_duration_seconds": 2.58,
    "total_recovery_duration_seconds": 11.28,
    "success": true
  }
}
\end{lstlisting}